\chapter{Launching}\label{ch:launching}

\section{Staffing}\label{sec:staffing}

Nella fase di planning si è stimato essere necessario avere 5 sviluppatori full-time per completare il progetto. Per questo motivo, il team che ha preso parte alla fase di planning corrisponde con il developer team: Asia, Giacomo, Linda (che ha anche ricoperto il ruolo di product owner), Nicolas e Nicolò.
Tutti i membri del team hanno già grande esperienza con progetti Scrum, per questo motivo non è stato necessario ingaggiare uno Scrum Master che guidi il team nell'applicazione di questa metodologia.

I membri del team sono ``persone fatte a T'': con questo termine si indica una persona con capacità di comprendere e affrontare uno spettro ampio di temi e problematiche -- rappresentato dalla barra orizzontale della T -- e allo stesso tempo con una forte competenza in una specifica materia -- rappresentata dalla barra verticale. Questo permette al team di avere una grande varietà di competenze e di poter affrontare qualsiasi problema che si presenti durante lo sviluppo del progetto. In seguito sono riportati i punti di forza di ciascuno sviluppatore:
\begin{itemize}
  \item Asia (\emph{accommodating}): ha grande esperienza nello sviluppo frontend, di interfacce utente con particolare attenzione alla user experience
  \item Giacomo (\emph{assimilating}): è esperto nella realizzazione di sistemi basati su una architettura funzionale e in linguaggi di programmazione funzionali 
  \item Linda (\emph{accommodating}): è esperta di Domain Driven Design e di modellazione di domini complessi
  \item Nicolas (\emph{diverging}): ha grande esperienza nell'utilizzare pratiche agili e strumenti automatizzati per garantire la qualità del codice, la scalabilità del sistema e la sicurezza delle operazioni
  \item Nicolò (\emph{converging}): è esperto nella progettazione, training e valutazione di sistemi di machine learning
\end{itemize}

\begin{figure}
  \centering
  \begin{tikzpicture}
    \draw[line width=30pt] (0,0) -- (0,4) node [midway, rotate=90, color=white] {conoscenze};  % linea verticale
    \draw[line width=30pt] (-2,4) -- (2,4) node[midway, color=white] {competenze}; % linea orizzontale
  \end{tikzpicture}
  \caption{Persona a T}
  \label{fig:croce}
\end{figure}

\todo[inline]{Decidere se fare o meno il Project Definition Statement}

\section{Definizione delle regole operative}
\subsection{Project meeting}
È stata definita la gamma di possibili meeting a cui il team dovrà partecipare, la loro durata e cadenza.

\subsubsection{Standup meeting}
Gli \emph{standup meeting} si svolgono con cadenza giornaliera e non impiegano più di 10 minuti. Questi meeting coinvolgono tutti i membri in piedi ogni mattina prima di iniziare la giornata.
Questi si svolgono in maniera informale e non hanno una scaletta esplicita: i membri del team di sviluppo discutono su ciò che stanno facendo e presentano gli eventuali problemi riscontrati. 
Nel caso in cui siano presenti problemi che impediscano di proseguire con lo sviluppo, il team si riunirà in un \emph{problem resolution meeting} per discutere più nel dettaglio del problema.

\subsubsection{Problem Resolution Meeting}
I problem resolution meeting vengono indetti in maniera straordinaria qualora durante uno standup meeting emerga un problema grave. Alla riunione partecipano il product owner, il \emph{problem owner} ed eventuali sviluppatori che possono essere affetti dal problema. Le riunioni si svolgono secondo una scaletta predefinita:
\begin{itemize}
  \item Vengono raccolte informazioni sul problema
  \item Si cercano possibili soluzioni tramite \emph{brainstorming}
  \item I partecipanti si consultano per decidere la soluzione migliore; nel caso in cui sia necessario potrebbero essere interpellati anche altri membri del team per chiedere un loro parere
  \item Si aggiorna lo \emph{sprint backlog} in modo da includere i task necessari a implementare la soluzione scelta al problema
\end{itemize}

\subsubsection{Sprint review e gestione del cambiamento di scope}
Il venerdì, al termine di ogni sprint, si svolge una riunione di \emph{sprint review} in cui il team presenta il lavoro svolto durante lo sprint. Alla riunione partecipano tutti i membri del team di sviluppo, il product owner, Raffaella (in quanto project sponsor) ed eventuali process owner lato cliente interessati ai risultati dello sprint.
Durante la riunione il team di sviluppo riceve feedback da parte del cliente mostrando le funzionalità implementate. Questo è anche un importante momento per mostrare eventuali \emph{mockup} per chiarire requisiti incerti.
Indicativamente queste riunioni impiegano da una a due ore a seconda del numero di user story implementate durante lo sprint.

Qualora dovessero emergere nuove richieste del committente o proposte di cambiamento di scope durante la fase di review la riunione potrebbe essere integrata con una fase in cui vengono discusse le nuove user story con il cliente, approvate e inserite nel product backlog secondo le metodologie già descritte nella fase di planning.

\subsubsection{Sprint retrospective}
Dopo ogni riunione di sprint review, il team di sviluppo si riunisce per un'ulteriore riunione di \emph{sprint retrospective}. Lo scopo di questa riunione è quello di analizzare il lavoro svolto durante lo sprint, cosa è andato bene, eventuali problematiche emerse e di discutere su come migliorare il processo di sviluppo. 
La durata di queste riunioni è indicativamente di 30 minuti.

\subsubsection{Sprint planning}
Si è fissata come data per le riunioni di \emph{sprint planning} il lunedì mattina all'inizio di ogni nuovo sprint.Durante queste riunioni il team di sviluppo si riunisce con il product owner per discutere delle user story da implementare durante lo sprint. Il ruolo del product owner è guidare il team nella scelta delle user story che forniscano il massimo valore possibile al cliente. Un fattore fondamentale nella scelta delle user story -- oltre al valore percepito dal cliente -- sono gli story points: infatti, si vuole che ogni sprint abbia un numero di story points adeguato alle risorse a disposizione e che lasci un margine di tempo libero per eventuali imprevisti.

Una volte scelte le user story, il team discute di come implementarle e di quali task sono necessari per completarle; una volta avvenuta la scomposizione, i task sono inseriti nello \emph{sprint backlog}. 

Successivamente, viene realizzato uno \emph{sprint network diagram}. Tale diagramma è analogo a un \emph{project network diagram} ma viene ripetuto ad ogni sprint per i suoi task. 
L'obiettivo di questo diagramma è comprendere le durate di ciascun task, le loro dipendenze e individuare il percorso critico. 
Questo diagramma ha un ruolo fondamentale nel monitoraggio dell'andamento dello sprint e permette di capire se eventuali problemi che emergono durante i daily meeting interessano il percorso critico comportando possibili ritardi nel completamento dello sprint. 

Infine, ogni programmatore si assegna i task che intende implementare durante lo sprint in base alle proprie competenze e alle priorità assegnate nello sprint network diagram.

Le riunioni di sprint planning impiegano circa 2 ore.

\subsubsection{Decision making style}
Date le dimensioni contenute del team di sviluppo e la sinergia fra i membri, si è optato per un approccio collaborativo al decision making. Ogni membro del team partecipa attivamente al processo decisionale esprimendo le proprie opinioni.
Nei casi in cui non si riesca a trovare una soluzione consensuale, il product owner ha il potere di prendere una decisione finale.

\subsubsection{Gestione delle comunicazioni}
Si è stabilito che le comunicazioni debbano avvenire secondo le seguenti modalità:
\begin{itemize}
  \item Email: per comunicare con il cliente in maniera asincrona; il punto di contatto lato team di sviluppo è Linda
  \item Slack: per comunicazioni interne al team di sviluppo
  \item Riunioni: preferibilmente in presenza, sia presso il caseificio Mambelli che la sede di Atedeg. In caso di impossibilità di svolgimento in presenza si è scelto di adottare Zoom come strumento per le riunioni in remoto
\end{itemize}
